\begin{helper}[Exact enumeration of one active support class]
Let an ETT with (n>0) noninitial rounds be fixed, let (q\in
\operatorname{Fin}(n)) satisfy (q+1=n), and fix the terminal support
(v_n\in V(T)).  Let (w) be a vertex whose
\noderef{ExtensionBaseSupportClass} (L(w)) is nonempty.  Assume that every
active PE support has the same \noderef{ETTFirstAppearanceIndex} as (v_n).
Then there are (p\geq0) and indices
[
                 i_0<i_1<\cdots<i_p
]
such that (w\in V(T_{q+1})), the range of this increasing sequence is
exactly (L(w)), and it is also exactly the set of all indices (r\leq q)
whose positive slot is a PE with exact chosen vertex (w).
\end{helper}
\begin{proof}
Choose (r_0\in L(w)).  By the definitions of
\noderef{ExtensionBaseSupportClass} and
\noderef{ExtensionBaseActiveIndices}, (r_0) is active and its chosen-vertex
field is both “some (w)” and “some (u)” for a vertex
(u\in V(T)).  Injectivity of the option constructor gives (u=w), so
(w\in V(T)), the slot (r_0+1) is a PE, and the common-index hypothesis
gives
[
 m(w)=m(v_n).                                                    \tag{1}
]

Now let (r\leq q) be any PE slot with chosen vertex (w).  The certificate
at (r+1), supplied by \noderef{TashkinovSeries}, must be the PE constructor
of \noderef{TashkinovSeriesStep}: its exact kind rules out the initial, SE,
and RE constructors.  In that PE constructor, the chosen vertex is the
vertex of its \noderef{MaximumDefectiveChoice}.  Its
\(\mathsf{vertexDefective}\) field, unfolded using
\noderef{DefectiveVertex}, is a conjunction whose membership clause puts
(w) in the current ladder tree (T_{r+1}).  By the
least-stage definition of \noderef{ETTFirstAppearanceIndex}, this implies
(m(w)\leq r+1).  Together with (1), the PE kind, the exact chosen-vertex
field, and (w\in V(T)), this is precisely membership of (r) in the active
set and hence in (L(w)).

Conversely, a member of (L(w)) is active and has chosen vertex (w) by
the two filter definitions.  It is therefore a PE.  Every index in
\(\operatorname{Fin}(n)\) is at most (q), because (q+1=n), so it satisfies
the displayed terminal bound.  Thus (L(w)) is exactly the claimed set of
terminal-bounded PE slots supported at (w).

The nonempty finite linearly ordered set (L(w)) has a unique increasing
enumeration by a type \(\operatorname{Fin}(p+1)\); take its cardinality minus
one for (p).  This gives the two asserted range equivalences.  Finally the
PE certificate at (r_0+1) puts (w) in (T_{r_0+1}).  At every later
certified slot, each constructor of \noderef{TashkinovSeriesStep} contains an
ordered initial-segment field.  Repeated application of
\noderef{TreeSequenceInitialSegmentVertexSetSubset} carries (w) through the
ladder trees up to (T_{q+1}), proving the remaining membership assertion.
\end{proof}
